\begin{tikzpicture}[
  font=\scriptsize,
  tanya/.style={draw, rounded corners, fill=orange!12, align=center,
    text width=2.55cm, minimum height=0.95cm, font=\scriptsize\bfseries},
  aturan/.style={draw, rounded corners, fill=blue!8, align=left,
    text width=2.55cm, inner sep=4pt, font=\tiny},
  induk/.style={draw, rounded corners, fill=praditagreen!15, align=center,
    text width=11.6cm, minimum height=0.6cm, font=\scriptsize\bfseries},
  hasil/.style={draw, rounded corners, fill=green!10, align=center,
    text width=11.6cm, minimum height=0.7cm, font=\scriptsize\bfseries},
  panah/.style={-{Stealth[length=2mm]}, thick}
]
  \node[induk] (mulai) at (4.5,1.3)
    {Satu \textit{workload} beserta jenis data yang diolahnya};

  \node[tanya] (t1) at (0,0) {Siapa penyelenggara sistemnya?};
  \node[tanya] (t2) at (3,0) {Ada data kesehatan atau rekam medis?};
  \node[tanya] (t3) at (6,0) {Ada data pribadi?};
  \node[tanya] (t4) at (9,0) {Ada aturan sektor atau kontrak lain?};

  \node[aturan, anchor=north] (a1) at (0,-0.95)
    {\textbf{PP 71/2019}\\
     Lingkup publik: di Indonesia, kecuali teknologinya tidak tersedia
     (Pasal 20).\\
     Lingkup privat: boleh di luar negeri bila pengawasan dan akses
     terjamin (Pasal 21).};
  \node[aturan, anchor=north] (a2) at (3,-0.95)
    {\textbf{PP 28/2024}\\
     Pusat data kesehatan di dalam wilayah Indonesia (Pasal 966); transfer dengan izin Menteri
     (Pasal 972).\\
     \textbf{Permenkes 24/2022}\\
     \textit{Cloud} tersertifikasi; cadangan di tempat berbeda (Pasal 20).\\
     Disimpan paling singkat 25 tahun (Pasal 39).};
  \node[aturan, anchor=north] (a3) at (6,-0.95)
    {\textbf{UU 27/2022 (PDP)}\\
     Data kesehatan termasuk data pribadi spesifik.\\
     Transfer ke luar negeri wajib memenuhi syarat (Pasal 56).\\
     Aturan pelaksana PP 33/2026 dipantau.};
  \node[aturan, anchor=north] (a4) at (9,-0.95)
    {\textbf{Aturan lain}\\
     Perjanjian kerja sama dengan BPJS Kesehatan, syarat mitra
     asuransi, atau kebijakan internal organisasi.};

  \node[hasil] (akhir) at (4.5,-4.85)
    {Lokasi data yang diizinkan dan daftar pertanyaan terbuka untuk
     \textit{workload} tersebut};

  \draw[panah] (mulai.south -| t1.north) -- (t1.north);
  \draw[panah] (t1.east) -- (t2.west);
  \draw[panah] (t2.east) -- (t3.west);
  \draw[panah] (t3.east) -- (t4.west);
  \draw[panah] (t1.south) -- node[right, font=\tiny] {selalu} (a1.north);
  \draw[panah] (t2.south) -- node[right, font=\tiny] {ya} (a2.north);
  \draw[panah] (t3.south) -- node[right, font=\tiny] {ya} (a3.north);
  \draw[panah] (t4.south) -- node[right, font=\tiny] {ya} (a4.north);
  \foreach \a in {a1, a2, a3, a4} {
    \draw[panah] (\a.south) -- (\a.south |- akhir.north);
  }
\end{tikzpicture}
